\chapter{Теоретическая часть}\label{ux442ux435ux43eux440ux435ux442ux438ux447ux435ux441ux43aux430ux44f-ux447ux430ux441ux442ux44c}

\section{Понятие генеративных моделей в машинном обучении}\label{ux43fux43eux43dux44fux442ux438ux435-ux433ux435ux43dux435ux440ux430ux442ux438ux432ux43dux44bux445-ux43cux43eux434ux435ux43bux435ux439-ux432-ux43cux430ux448ux438ux43dux43dux43eux43c-ux43eux431ux443ux447ux435ux43dux438ux438}

Машинное обучение представляет собой раздел искусственного интеллекта, посвящённый разработке алгоритмов, способных извлекать закономерности из данных и применять их к новым, ранее не встречавшимся примерам {[}19, 20{]}. В зависимости от характера решаемой задачи методы машинного обучения принято делить на несколько крупных классов: обучение с учителем, обучение без учителя, обучение с подкреплением и генеративное обучение. Каждый из этих классов имеет собственные принципы построения моделей и собственные критерии качества.

При обучении с учителем модель получает примеры с правильными ответами. Обучение без учителя предполагает самостоятельный поиск структур и кластеров в данных, а обучение с подкреплением строится на взаимодействии со средой и получении награды или штрафа.

Генеративное обучение занимает в этой классификации особое место. Его задача -- не предсказать ответ по входным данным, а научиться порождать новые данные, похожие на те, что были в обучающей выборке {[}1, 13{]}. Иными словами, модель должна усвоить вероятностное распределение данных и затем уметь извлекать из него новые примеры. Если модель обучена на фотографиях кошек, она должна уметь генерировать новые фотографии кошек, которых никогда не было, но которые выглядят правдоподобно.

Классический пример генеративной модели -- GAN, предложенная И. Гудфеллоу и соавторами в 2014 году {[}13{]}. Она объединяет генератор, создающий образцы, и дискриминатор, отличающий их от реальных. Соревнование сетей постепенно повышает качество генерации. GAN широко применяются для изображений, видео и звука, но менее удобны для текста из-за его дискретной природы.

Языковые модели оценивают вероятность последовательностей символов или слов и порождают новые последовательности {[}5, 17{]}. Развитие этого направления прошло путь от частотных n-граммных моделей к рекуррентным сетям и трансформерам {[}5, 11{]}. Большие языковые модели, например GPT, создают связный текст, но требуют значительных вычислительных ресурсов и больших обучающих корпусов {[}18{]}.

В данной работе рассматриваются посимвольные языковые модели для генерации коротких последовательностей -- названий. Они не требуют токенизации, могут обучаться на сравнительно небольших датасетах и запускаться на обычном компьютере. Это делает их подходящими для узкой предметной области растений и продуктов.

Обучение и инференс -- разные этапы работы генеративной модели. При обучении параметры корректируются на примерах из датасета; при инференсе уже обученная модель создаёт новые последовательности. В работе обучение выполняется локально, а генерация доступна через приложение.

\section{Задача генерации текстовых последовательностей}\label{ux437ux430ux434ux430ux447ux430-ux433ux435ux43dux435ux440ux430ux446ux438ux438-ux442ux435ux43aux441ux442ux43eux432ux44bux445-ux43fux43eux441ux43bux435ux434ux43eux432ux430ux442ux435ux43bux44cux43dux43eux441ux442ux435ux439}

Генерация текста заключается в построении модели, которая по обучающим примерам создаёт новые статистически похожие последовательности {[}4, 5{]}. Для названий растений и продуктов важно, чтобы результат был правдоподобным, соответствовал выбранной категории и не повторял примеры корпуса.

Один из первых подходов -- n-граммные модели, оценивающие частоты соседних символов или слов {[}5, 10{]}. Модель порядка n учитывает предыдущие n−1 элементов и на их основе выбирает следующий. Например, биграммная модель использует пары соседних символов, поэтому её работу легко понять и проверить.

Недостаток n-грамм -- слабая работа с редкими сочетаниями: если нужный переход не встречался, оценить его вероятность трудно. Сглаживание Кнесера -- Нея перераспределяет вероятности и даёт небольшую вероятность редким сочетаниям {[}10{]}.

При увеличении n число возможных контекстов быстро растёт, а данных для их оценки не хватает. Поэтому на практике используют небольшие порядки, что ограничивает учёт дальнего контекста.

Нейросетевые языковые модели заменили таблицы частот обучаемыми векторными представлениями и сетями, которые выявляют закономерности сочетания символов и слов {[}5, 15, 16{]}. Модель Бенджио и соавторов 2003 года стала одним из ранних примеров такого подхода {[}5{]}.

Посимвольные модели рассматривают название как последовательность отдельных знаков. Такой подход не требует словаря готовых слов и позволяет получать неологизмы {[}7, 16{]}. Его недостаток -- большее число шагов по сравнению со словесной моделью, что увеличивает вычислительную нагрузку.

Генерация носит вероятностный характер: на каждом шаге следующий символ выбирается из распределения вероятностей. Благодаря этому одна модель создаёт разные варианты. Степень случайности регулируется параметрами сэмплирования.

Качество генерации оценивают количественно -- по перплексии, новизне и разнообразию -- и качественно, с помощью экспертной оценки. Поскольку единственного правильного ответа нет, метрики рассматриваются совместно; подробнее они описаны в разделе 1.5.

\section{Рекуррентные нейронные сети и архитектура LSTM}\label{ux440ux435ux43aux443ux440ux440ux435ux43dux442ux43dux44bux435-ux43dux435ux439ux440ux43eux43dux43dux44bux435-ux441ux435ux442ux438-ux438-ux430ux440ux445ux438ux442ux435ux43aux442ux443ux440ux430-lstm}

Рекуррентные нейронные сети (RNN) предназначены для обработки последовательностей {[}1, 6{]}. Их скрытое состояние передаётся между шагами, поэтому при обработке текущего символа сеть учитывает предшествующий контекст. В отличие от полносвязных сетей, элементы последовательности не рассматриваются независимо.

На каждом шаге RNN получает текущий символ и предыдущее скрытое состояние, а возвращает новое состояние и прогноз следующего символа. В скрытом состоянии сохраняется информация о ранее обработанной части последовательности.

Простые RNN подвержены затуханию и взрыву градиента {[}6{]}. При обратном распространении ошибки через много временных шагов градиенты могут становиться слишком малыми или слишком большими. В результате сеть забывает дальний контекст либо обучается нестабильно.

Решением проблемы стала архитектура LSTM (Long Short-Term Memory, «долгая краткосрочная память»), предложенная в 1997 году З. Хохрайтером и Ю. Шмидхубером {[}6{]}. LSTM сохраняет общую идею рекуррентной обработки, но вводит более сложную внутреннюю структуру ячейки, способную эффективно сохранять и передавать информацию на длинные дистанции.

Ключевым элементом LSTM является наличие двух состояний: скрытого состояния (как в обычной RNN) и так называемого состояния ячейки (cell state). Состояние ячейки представляет собой долговременную память, которая может беспрепятственно переноситься через многие шаги с минимальными изменениями. Управление этой памятью осуществляется с помощью трёх специальных механизмов, называемых «вентилями» (gates): вентиля забывания, входного вентиля и выходного вентиля.

Вентиль забывания определяет, какую информацию из предыдущего состояния сохранить. Входной вентиль регулирует запись новых сведений, а выходной -- передачу части состояния ячейки в скрытое состояние для прогноза следующего символа.

Благодаря такой структуре LSTM может сохранять важные сведения на протяжении многих шагов и отбрасывать несущественные. Архитектура применялась для языкового моделирования {[}15{]}, генерации текста {[}7, 16{]}, машинного перевода и распознавания речи.

В контексте генерации названий LSTM работает следующим образом. На вход модели подаётся начало названия (или специальный маркер начала). На каждом шаге модель предсказывает распределение вероятностей по всем возможным символам и выбирает один из них. Выбранный символ добавляется к последовательности и подаётся на вход на следующем шаге. Процесс повторяется до тех пор, пока не будет сгенерирован специальный маркер конца или не будет достигнута максимальная длина.

Для условной генерации в модель дополнительно передаются категория названия и требуемое количество слов. Их векторные представления объединяются с символьными эмбеддингами и подаются на вход LSTM. Благодаря этому модель учитывает запрос пользователя, например генерирует названия именно для ягод или зерновых.

Кроме того, при обучении LSTM применяются стандартные методы регуляризации -- в частности, dropout {[}14{]}, который случайным образом «выключает» часть нейронов во время обучения, что помогает предотвратить переобучение. Используется также оптимизатор Adam {[}12{]} -- адаптивный метод градиентного спуска, хорошо зарекомендовавший себя в задачах обучения нейронных сетей.

Хотя трансформеры получили широкое распространение {[}11{]}, LSTM остаётся практичным выбором для небольших задач. Такие модели компактны, обучаются на CPU и требуют сравнительно небольшого датасета, что соответствует условиям настоящей работы.

Таким образом, LSTM сочетает последовательную обработку с механизмом долговременной памяти и подходит для генерации коротких названий. Далее рассматривается альтернативная архитектура -- трансформер.

\section{Трансформеры и механизм внимания}\label{ux442ux440ux430ux43dux441ux444ux43eux440ux43cux435ux440ux44b-ux438-ux43cux435ux445ux430ux43dux438ux437ux43c-ux432ux43dux438ux43cux430ux43dux438ux44f}

Архитектура LSTM, рассмотренная в предыдущем разделе, долгое время оставалась основной для задач обработки последовательностей. Однако у неё есть принципиальное ограничение: обработка последовательности происходит строго последовательно, шаг за шагом. Это затрудняет распараллеливание вычислений и ограничивает длину контекста, который модель может эффективно учитывать. Решением этой проблемы стала архитектура трансформера, предложенная в 2017 году группой исследователей из Google под руководством А. Васвани {[}11{]}.

Ключевой идеей трансформера является механизм внимания (attention mechanism) -- специальный способ обработки последовательности, при котором каждый элемент может напрямую «взаимодействовать» с любым другим элементом, независимо от их взаимного расположения. В отличие от рекуррентных сетей, где информация передаётся последовательно через скрытые состояния, в трансформере каждый элемент «видит» всю последовательность целиком и сам решает, на какие элементы обратить внимание.

Механизм внимания удобно представить следующим образом. Для каждого элемента последовательности вычисляются три векторных представления: запрос (query), ключ (key) и значение (value). Запрос описывает, что именно ищет текущий элемент; ключ -- что может предложить каждый элемент; значение -- какую информацию несёт каждый элемент. Затем для текущего элемента вычисляется степень его соответствия каждому элементу последовательности -- насколько его запрос «подходит» к ключам других элементов. По этим степеням формируются веса, с которыми значения других элементов суммируются и добавляются к представлению текущего элемента. В результате каждый элемент получает «сжатую» информацию обо всей последовательности, взвешенную по релевантности.

Такая схема получила название self-attention (самовнимание), поскольку запросы, ключи и значения вычисляются из одной и той же последовательности. Она позволяет модели гибко учитывать контекст: например, при обработке названия «Пшеница Озимая» модель может «обратить внимание» на связь между первым и вторым словом, чтобы правильно определить род прилагательного.

Одним из важных усовершенствований механизма внимания является многоголовое внимание (multi-head attention). Идея в том, что вместо одной «головы» внимания используется несколько параллельных «голов», каждая из которых фокусируется на своём аспекте последовательности. Например, одна голова может отслеживать синтаксические связи, другая -- фонетические паттерны, третья -- семантические. Результаты всех голов затем объединяются. Это делает модель более выразительной и устойчивой.

Полноценная архитектура трансформера состоит из энкодера и декодера, каждый из которых содержит несколько слоёв. Слой энкодера включает механизм многоголового self-attention и полносвязную нейронную сеть с нелинейной активацией; слой декодера дополнительно содержит механизм attention между энкодером и декодером. В задачах генерации текста (в том числе генерации названий) чаще используется декодер-only вариант, в котором модель просто предсказывает следующий символ по предыдущим, но делает это с помощью механизма внимания вместо рекуррентности.

По сравнению с LSTM, трансформер обладает несколькими существенными преимуществами. Во-первых, он легко распараллеливается: все позиции обрабатываются одновременно, что позволяет эффективно использовать графические ускорители и обучать очень большие модели. Во-вторых, он лучше справляется с длинными зависимостями: информация между удалёнными элементами передаётся напрямую, без искажений через промежуточные состояния. В-третьих, его архитектура более гибкая: добавляя слои и головы внимания, можно легко масштабировать модель.

Вместе с тем трансформер имеет и недостатки, особенно важные для прикладных задач малого масштаба. Главный из них высокие требования к объёму данных. Из-за большого числа параметров трансформеру нужно гораздо больше примеров для обучения, чем LSTM. На небольших датасетах (несколько сотен или тысяч примеров) трансформер склонен к переобучению и уступает более простым архитектурам. Второй недостаток -- вычислительная стоимость: обучение и инференс трансформера заметно дороже, чем у LSTM, даже при сопоставимом числе параметров.

\textbf{Таблица 1.1 -- Сравнение архитектур LSTM и Трансформер}

\begin{longtable}[]{@{}
  >{\raggedright\arraybackslash}p{(\columnwidth - 4\tabcolsep) * \real{0.3243}}
  >{\raggedright\arraybackslash}p{(\columnwidth - 4\tabcolsep) * \real{0.3243}}
  >{\raggedright\arraybackslash}p{(\columnwidth - 4\tabcolsep) * \real{0.3243}}@{}}
\toprule\noalign{}
\begin{minipage}[b]{\linewidth}\raggedright
\textbf{Критерий}
\end{minipage} & \begin{minipage}[b]{\linewidth}\raggedright
\textbf{LSTM}
\end{minipage} & \begin{minipage}[b]{\linewidth}\raggedright
\textbf{Трансформер}
\end{minipage} \\
\midrule\noalign{}
\endhead
\bottomrule\noalign{}
\endlastfoot
Принцип обработки последовательности & Последовательный (шаг за шагом) & Параллельный (все позиции одновременно) \\
Учёт долгосрочных зависимостей & Через скрытое состояние и вентили & Через механизм внимания напрямую \\
Требования к объёму данных & Умеренные (сотни--тысячи примеров) & Высокие (десятки тысяч примеров) \\
Вычислительная стоимость обучения & Низкая (работает на CPU) & Высокая (желателен GPU) \\
Склонность к переобучению на малых данных & Низкая & Высокая \\
Скорость инференса & Высокая & Средняя \\
Пригодность для малых датасетов & Высокая & Низкая \\
\end{longtable}

Именно поэтому в настоящей курсовой работе основным выбором стала LSTM, а трансформер использовался в качестве альтернативной архитектуры для сравнения. На датасете из нескольких тысяч названий LSTM показала более стабильные результаты и лучше обобщила на новые примеры, тогда как трансформер продемонстрировал склонность к переобучению. Этот вывод согласуется с общей практикой применения нейросетевых моделей: выбор архитектуры определяется не только её теоретической мощностью, но и объёмом доступных данных.

Следует также отметить, что для реальных приложений генерации названий малого масштаба не требуется мощная языковая модель уровня GPT. Как показала практика, компактная LSTM, обученная на нескольких тысячах примеров, способна генерировать качественные и осмысленные названия при минимальных вычислительных затратах. Трансформеры же показывают свои преимущества в задачах с большими объёмами данных и длинными контекстами -- в тех сценариях, где их вычислительная стоимость оправдана.

Таким образом, механизм внимания и архитектура трансформера представляют собой важный этап в развитии методов обработки последовательностей. Они расширяют возможности моделирования долгосрочных зависимостей и обеспечивают высокую производительность на больших данных. Однако для задач малого масштаба, к которым относится генерация названий растений и продуктов, LSTM остаётся более практичным выбором. В следующем разделе будут рассмотрены методы оценки качества сгенерированных названий -- независимо от того, какой архитектурой они получены.

\section{Методы оценки качества сгенерированных названий}\label{ux43cux435ux442ux43eux434ux44b-ux43eux446ux435ux43dux43aux438-ux43aux430ux447ux435ux441ux442ux432ux430-ux441ux433ux435ux43dux435ux440ux438ux440ux43eux432ux430ux43dux43dux44bux445-ux43dux430ux437ux432ux430ux43dux438ux439}

Оценка генеративных моделей сложнее оценки классификаторов: у сгенерированного названия нет единственного правильного ответа. Поэтому для сравнения моделей используются несколько взаимодополняющих метрик: перплексия, стилевое соответствие, новизна, разнообразие и человеческая оценка.

Перплексия (perplexity) показывает, насколько уверенно языковая модель предсказывает последовательность {[}4, 15{]}. Чем ниже значение, тем лучше модель описывает данные. При равномерном случайном выборе символов перплексия близка к размеру словаря, а при идеальном предсказании равна единице.

Перплексия является относительной метрикой: корректно сравнивать её можно только для моделей с одинаковым словарём и датасетом. Она характеризует предсказание последовательностей, но не оценивает эстетическое качество названий.

Стилевое соответствие (style adherence) показывает, насколько результат отвечает заданной категории. Для этого используется классификатор, обученный распознавать категории названий. Оценкой служит вероятность класса, соответствующего запросу пользователя.

Такой способ не требует ручной разметки каждого нового результата: классификатор применяет закономерности, усвоенные на исходном датасете.

Новизна (novelty) показывает, насколько сгенерированное название отличается от обучающих примеров. Для её оценки можно сравнивать символьные n-граммы нового названия с n-граммами корпуса. Чем меньше сходство с существующими названиями, тем выше новизна. При этом полное отсутствие общих фрагментов не всегда означает осмысленность результата, поэтому метрику нельзя использовать отдельно.

Разнообразие (diversity) характеризует различия между названиями, созданными в одной серии. Если модель постоянно повторяет одинаковые или почти одинаковые варианты, разнообразие низкое. Для оценки применяются доля уникальных результатов и расстояния между строками.

Человеческая оценка (human evaluation) остаётся важным способом проверки, хотя требует времени. Несколько респондентов могут оценить естественность, соответствие категории и удобство использования названий.

В настоящей работе такая оценка проводилась в ограниченном объёме: респонденты сравнивали варианты и отмечали наиболее удачные. Результаты использовались как дополнительная качественная проверка.

Для сопоставления также использовались базовые модели -- марковская цепь и n-граммный генератор. Сравнение с ними показывает, даёт ли нейросетевая архитектура практическое преимущество.

Таким образом, качество оценивалось по совокупности языковых, стилевых и сравнительных показателей, а также с помощью человеческой оценки.

\textbf{Таблица 1.2 -- Метрики оценки качества генеративных моделей}

\begin{longtable}[]{@{}
  >{\raggedright\arraybackslash}p{(\columnwidth - 6\tabcolsep) * \real{0.2000}}
  >{\raggedright\arraybackslash}p{(\columnwidth - 6\tabcolsep) * \real{0.2533}}
  >{\raggedright\arraybackslash}p{(\columnwidth - 6\tabcolsep) * \real{0.2400}}
  >{\raggedright\arraybackslash}p{(\columnwidth - 6\tabcolsep) * \real{0.2800}}@{}}
\toprule\noalign{}
\begin{minipage}[b]{\linewidth}\raggedright
\textbf{Метрика}
\end{minipage} & \begin{minipage}[b]{\linewidth}\raggedright
\textbf{Что измеряет}
\end{minipage} & \begin{minipage}[b]{\linewidth}\raggedright
\textbf{Диапазон значений}
\end{minipage} & \begin{minipage}[b]{\linewidth}\raggedright
\textbf{Интерпретация}
\end{minipage} \\
\midrule\noalign{}
\endhead
\bottomrule\noalign{}
\endlastfoot
Перплексия & Уверенность модели в предсказаниях & От 1 до размера словаря & Чем ниже, тем лучше \\
Стилевое соответствие & Доля названий, соответствующих категории & От 0 до 1 & Чем выше, тем лучше \\
Новизна & Доля названий, отсутствующих в датасете & От 0 до 1 & Чем выше, тем лучше \\
Разнообразие & Энтропия распределения символов & От 0 до log₂(размер словаря) & Чем выше, тем лучше \\
Человеческая оценка & Субъективное восприятие респондентов & От 1 до 5 & Чем выше, тем лучше \\
\end{longtable}

\section{Обзор существующих генераторов названий}\label{ux43eux431ux437ux43eux440-ux441ux443ux449ux435ux441ux442ux432ux443ux44eux449ux438ux445-ux433ux435ux43dux435ux440ux430ux442ux43eux440ux43eux432-ux43dux430ux437ux432ux430ux43dux438ux439}

Задача генерации названий имеет давнюю историю и большое число практических реализаций. Можно выделить несколько классов существующих инструментов: от простых генераторов на основе шаблонов до современных нейросетевых сервисов.

Шаблонные генераторы используют заранее подготовленные списки слогов или частей слов. Они просты и быстры, но дают ограниченное разнообразие и слабо учитывают особенности конкретной предметной области.

Марковские генераторы строят последовательности на основе вероятностей перехода между символами. Они способны создавать новые комбинации, однако качество зависит от объёма корпуса и порядка модели; дальние зависимости обычно не учитываются.

N-граммные генераторы со сглаживанием расширяют этот подход за счёт учёта более длинных фрагментов и ненулевых вероятностей редких сочетаний. Тем не менее их возможности ограничены размером обучающего корпуса.

Нейросетевые генераторы используют LSTM или GRU и обучаются предсказывать следующий символ или слово. Они лучше обобщают последовательности и могут учитывать более длинный контекст, но требуют подготовки данных и настройки параметров.

Большие языковые модели (LLM), включая GPT, Claude и LLaMA, способны генерировать связный текст по запросу. Однако для узкой задачи нейминга их использование может быть избыточным: требуется доступ к крупной модели, а результат не всегда соответствует заданным ограничениям без дополнительной фильтрации.

Для игр существуют генераторы фэнтезийных имён и топонимов. Обычно они ориентированы на общий художественный стиль, а не на предметные категории сельского хозяйства и продуктов питания.

Специализированных генераторов названий растений и продуктов немного. Разрабатываемая система ориентирована именно на эту область и позволяет управлять категорией и длиной результата.
